% 5.5 表示域与基本域
\section{表示域与基本域}

在表 5.7（以及 Miller and Love (1967) 的表）中，空间群表示只对 230 个空间群中每一个的布里渊区仅仅一个基本域（见定义 3.3.3）内的特殊对称点与对称线列出。回顾基本域 $\Omega$ 具有性质 $(\sum_{R}R\Omega)$ 等于整个布里渊区，其中算符 $R$ 取遍相应晶系全对称点群 $\mathbf{P}$ 的元素。例如对立方晶体，基本域只是布里渊区的四十八分之一。由于在讨论物理问题时可能会用到它们，为了完整性，必须考虑：给定仅仅一个基本域内的小表示之后，如何确定与布里渊区中其他点或对称线相联系的小表示。回顾表示域 $\Phi$（见定义 3.7.1）具有性质 $(\sum_{S}S\Phi)$ 等于整个布里渊区，其中算符 $S$ 取遍所考虑空间群 $\mathbf{G}$ 的同形点群 $\mathbf{F}$ 的元素。例如对空间群 $Pa3\ (T^{5}_{h})$，同形点群是 $m3\ (T_{h})$，含 24 个对称操作，故 $\Phi$ 是布里渊区的二十四分之一；而基本域 $\Omega$ 由全对称立方点群 $m3m\ (O_{h})$（含 48 个对称操作）决定，$\Omega$ 是布里渊区的四十八分之一。所以此情形表示域的大小是基本域的两倍。一般地，由于 $\mathbf{F}$ 或者等于 $\mathbf{P}$ 或者是 $\mathbf{P}$ 的真子群，表示域 $\Phi$ 要么与 $\Omega$ 相同，要么是若干个与 $\Omega$ 全等的体积分之和。

于是有两种情形要考虑：第一种是表示域与基本域一样大；第二种是表示域更大。若表示域 $\Phi$ 与基本域 $\Omega$ 一样大，即 $\mathbf{F} = \mathbf{P}$，则任何在基本域之外的波矢 $\mathbf{k}_{B}$ 都通过
\begin{equation*}
\mathbf{k}_{B} \equiv R\mathbf{k}_{A}
\tag{5.5.1}
\end{equation*}
与基本域内的某个波矢 $\mathbf{k}_{A}$ 相联系，其中 $R \in \mathbf{P}$。由于 $\mathbf{P} = \mathbf{F}$，这意味着此情形下总存在 $\mathbf{k}_{B}$ 的星中某个位于布里渊区基本域内的波矢 $\mathbf{k}_{A}$。于是与波矢为 $\mathbf{k}_{B}$ 的粒子或准粒子相联系的任何物理可观察量，将与波矢为 $\mathbf{k}_{A}$ 的粒子或准粒子所拥有的值相同。因此，若表示域 $\Phi$ 与基本域 $\Omega$ 一样大，研究基本域之外的任何波矢都不会带来新的物理见识。很可能正因为实际用于能带结构计算或声子色散关系研究的大多数空间群都属于这一类，才会产生"只需考虑布里渊区基本域内的波矢总是足够的"这一错误假设。

第二种情形，表示域 $\Phi$ 的体积是基本域 $\Omega$ 体积的某个小整数倍。但现在，对 $\Phi$ 中而不在 $\Omega$ 中的矢量 $\mathbf{k}_{B}$，不存在操作 $R \in \mathbf{F}$ 通过式 (5.5.1) 把 $\mathbf{k}_{B}$ 与基本域内的任何矢量 $\mathbf{k}_{A}$ 联系起来。此时满足式 (5.5.1) 的操作 $R$ 属于 $\mathbf{P}$ 但不属于 $\mathbf{F}$。然而仍然成立的是：从物理观点看所有显著不同的波矢，都可以通过考虑表示域 $\Phi$ 中的全部波矢而得到。例如，在任何情形下，能带结构计算的能级本征值曲面 $E(\mathbf{k}) = E$ 都将具有同形点群 $\mathbf{F}$ 的点群对称性。表示域之外的波矢仍然无需考虑，但必须能够确定位于 $\Phi$ 之内、$\Omega$ 之外的波矢 $\mathbf{k}_{B}$ 的空间群表示——它们不能直接从表 5.7 得到。为此我们首先注意到如下定理。

\noindent\textbf{定理 5.5.1}\quad 若 $\mathbf{F}$ 是给定晶系的一个点群，而 $\mathbf{P}$ 是该晶系的全对称点群，则 $\mathbf{F}$ 是 $\mathbf{P}$ 的不变子群。

证明最简单地是通过枚举与观察完成。我们前面已指出每个非平凡空间群都有指数为 2 或 3 的不变子群（Zak (1960)），由此可证：存在 24 个空间群（下述 (5.5.2)），它们虽然彼此同构可区分，却不能由表 5.7 直接给出基本域之外的表示。这 24 个空间群是
\begin{equation*}
\left.
\begin{array}{l}
\left\{P4_{1}(C^{2}_{4}),\ P4_{3}(C^{4}_{4})\right\},\ \left\{P3_{1}(C^{2}_{3}),\ P3_{2}(C^{3}_{3})\right\},\\
\left\{P3_{1}12(D^{3}_{3}),\ P3_{2}12(D^{5}_{3})\right\},\ \left\{P3_{1}21(D^{4}_{3}),\ P3_{2}21(D^{6}_{3})\right\},\\
\left\{P6_{1}(C^{2}_{6}),\ P6_{5}(C^{3}_{6})\right\},\ \left\{P6_{2}(C^{4}_{6}),\ P6_{4}(C^{5}_{6})\right\},\\
\left\{P4_{1}22(D^{3}_{4}),\ P4_{3}22(D^{7}_{4})\right\},\ \left\{P4_{1}2_{1}2(D^{4}_{4}),\ P4_{3}2_{1}2(D^{8}_{4})\right\},\\
P6_{1}22(D^{2}_{6}),\ P6_{5}22(D^{3}_{6}),\ P6_{2}22(D^{4}_{6}),\ P6_{4}22(D^{5}_{6}),\\
P2_{1}3(T^{4}),\\
Pa3(T^{6}_{h}),\\
P4_{3}32(O^{6})\ \text{与}\ P4_{1}32(O^{7}).
\end{array}
\right\}
\tag{5.5.2}
\end{equation*}
这 24 个空间群组成了 11 对同构但晶体学上可区分的群，连同两个例外群 $P2_{1}3\ (T^{4})$ 与 $Pa3\ (T^{6}_{h})$。

设 $\mathbf{k}_{B}$ 是 $\Phi$ 中而不在 $\Omega$ 中的波矢。则存在基本域 $\Omega$ 中的波矢 $\mathbf{k}_{A}$ 与一个点群操作 $R$——它不属于 $\mathbf{G}$ 的同形点群 $\mathbf{F}$，但属于相应晶系的全对称点群 $\mathbf{P}$——使得
\begin{equation*}
\mathbf{k}_{B} \equiv R\mathbf{k}_{A}.
\tag{5.5.1}
\end{equation*}
由于 $\mathbf{k}_{A}$ 在 $\Omega$ 中，小群 $\mathbf{G}^{\mathbf{k}_{A}}$ 及其小表示可得，且可从表 5.7 辨认。

\noindent\textbf{定理 5.5.2}\quad 若 $\overline{\mathbf{G}}^{\mathbf{k}_{A}}$ 是 $\mathbf{k}_{A}$ 的小陪群，则 $\mathbf{k}_{B}$ 的小陪群由 $\overline{\mathbf{G}}^{\mathbf{k}_{B}} = R\overline{\mathbf{G}}^{\mathbf{k}_{A}}R^{-1}$ 给出。

该定理可以如下证明。设 $S \in \overline{\mathbf{G}}^{\mathbf{k}_{A}}$，则
\begin{equation*}
S\mathbf{k}_{A} \equiv \mathbf{k}_{A},
\tag{5.5.3}
\end{equation*}
由式 (5.5.1) 可写为
\begin{equation*}
SR^{-1}\mathbf{k}_{B} \equiv R^{-1}\mathbf{k}_{B}
\tag{5.5.4}
\end{equation*}
因此
\begin{equation*}
(RSR^{-1})\mathbf{k}_{B} \equiv \mathbf{k}_{B}.
\tag{5.5.5}
\end{equation*}
现在 $S \in \overline{\mathbf{G}}^{\mathbf{k}_{A}} \subset \mathbf{F}$ 且 $R \in \mathbf{P}$。由于 $\mathbf{F}$ 在 $\mathbf{P}$ 中不变（见定理 5.5.1），可知 $RSR^{-1} \in \mathbf{F}$。这一事实与式 (5.5.5) 一起蕴含 $RSR^{-1} \in \overline{\mathbf{G}}^{\mathbf{k}_{B}}$。由于这对一切 $S \in \overline{\mathbf{G}}^{\mathbf{k}_{A}}$ 成立，故 $R\overline{\mathbf{G}}^{\mathbf{k}_{A}}R^{-1} \subseteq \overline{\mathbf{G}}^{\mathbf{k}_{B}}$。从任一元素 $T \in \overline{\mathbf{G}}^{\mathbf{k}_{B}}$ 出发，同样可以证明 $R^{-1}TR \in \overline{\mathbf{G}}^{\mathbf{k}_{A}}$，从而 $R^{-1}\overline{\mathbf{G}}^{\mathbf{k}_{B}}R \subseteq \overline{\mathbf{G}}^{\mathbf{k}_{A}}$，即 $\overline{\mathbf{G}}^{\mathbf{k}_{B}} \subseteq R\overline{\mathbf{G}}^{\mathbf{k}_{A}}R^{-1}$。于是反向包含也成立，因此
\begin{equation*}
\overline{\mathbf{G}}^{\mathbf{k}_{B}} = R\overline{\mathbf{G}}^{\mathbf{k}_{A}}R^{-1}.
\tag{5.5.6}
\end{equation*}
所以小陪群 $\overline{\mathbf{G}}^{\mathbf{k}_{B}}$ 与小陪群 $\overline{\mathbf{G}}^{\mathbf{k}_{A}}$ 共轭（相对于 $\mathbf{P}$）。$\overline{\mathbf{G}}^{\mathbf{k}_{B}}$ 与 $\overline{\mathbf{G}}^{\mathbf{k}_{A}}$ 之间显然有同构，因为若 $S_{i}, S_{j} \in \overline{\mathbf{G}}^{\mathbf{k}_{A}}$ 且
\begin{equation*}
S_{i}S_{j} = S_{k},
\tag{5.5.7}
\end{equation*}
则记 $T_{i} = RS_{i}R^{-1}$ 等，可得
\begin{equation*}
T_{i}T_{j} = RS_{i}R^{-1}RS_{j}R^{-1} = RS_{i}S_{j}R^{-1} = RS_{k}R^{-1} = T_{k}.
\tag{5.5.8}
\end{equation*}
很多时候小陪群 $\overline{\mathbf{G}}^{\mathbf{k}_{B}}$ 与 $\overline{\mathbf{G}}^{\mathbf{k}_{A}}$ 会重合。

现在的问题是：小群 $\mathbf{G}^{\mathbf{k}_{B}}$ 是否共轭于（相对于母空间群 $\mathbf{G}$ 的）小群 $\mathbf{G}^{\mathbf{k}_{A}}$。设 $\mathbf{D}^{\mathbf{k}_{A}}$ 是 $\mathbf{G}^{\mathbf{k}_{A}}$ 的一个允许的小表示（基 $\phi$），则可以证明（细节见原文）：$\mathbf{G}^{\mathbf{k}_{B}}$ 的表示由 $\phi' = \{R \mid \mathbf{v}\}\phi$ 给出基、按 $\mathbf{D}^{\mathbf{k}_{B}}(\{R \mid \mathbf{v}\}^{-1}g\{R \mid \mathbf{v}\}) = \mathbf{D}^{\mathbf{k}_{A}}(g)$ 相关联；即使 $\mathbf{G}^{\mathbf{k}_{B}}$ 与 $\mathbf{G}^{\mathbf{k}_{A}}$ 不同构，这一构造仍然给出 $\mathbf{G}^{\mathbf{k}_{B}}$ 的允许表示。

因此一般程序是：对 $\Phi \setminus \Omega$ 中的每个点或线 $\mathbf{k}_{B}$，找 $\mathbf{P}$ 中把 $\Omega$ 中的 $\mathbf{k}_{A}$ 变到 $\mathbf{k}_{B}$ 的操作 $R$；然后 $\overline{\mathbf{G}}^{\mathbf{k}_{B}} = R\overline{\mathbf{G}}^{\mathbf{k}_{A}}R^{-1}$ 给出新的小陪群，其表示由 $\mathbf{G}^{\mathbf{k}_{A}}$ 的表示经共轭得到。

作为示例，表 5.9 列出空间群 $F4_{1}32\ (O^{4})$ 的 $\sigma_{db}\Omega$ 中的对称点与对称线：基本域 $\Omega$ 中的 $\mathbf{k}_{A}$ 与 $\Phi \setminus \Omega$ 中的 $\mathbf{k}_{B} = \sigma_{db}\mathbf{k}_{A}$ 的对应。

\begin{table}[H]
\centering
\caption*{\textbf{表 5.9}\ $F4_{1}32\ (O^{4})$ 的 $\sigma_{db}\Omega$ 中的对称点与对称线}
\begin{tabular}{@{}lcll@{}}
\toprule
 & $\mathbf{k}_{A}$ & \multicolumn{2}{c}{$\mathbf{k}_{B} = \sigma_{db}\mathbf{k}_{A}$}\\
\midrule
$\Gamma$ & $(0, 0, 0)$ & $\Gamma$ & $(0, 0, 0)$\\
$X$ & $(\frac{1}{2}, 0, \frac{1}{2})$ & $X'$ & $(0, \frac{1}{2}, \frac{1}{2})$\\
$L$ & $(\frac{1}{2}, \frac{1}{2}, \frac{1}{2})$ & $L$ & $(\frac{1}{2}, \frac{1}{2}, \frac{1}{2})$\\
$W$ & $(\frac{1}{2}, \frac{1}{4}, \frac{3}{4})$ & $W'$ & $(\frac{1}{4}, \frac{1}{2}, \frac{3}{4})$\\
$\Delta$ & $(\alpha, 0, \alpha)$ & $\Delta'$ & $(0, \alpha, \alpha)$\\
$\Lambda$ & $(\alpha, \alpha, \alpha)$ & $\Lambda$ & $(\alpha, \alpha, \alpha)$\\
$\Sigma$ & $(\alpha, \alpha, 2\alpha)$ & $\Sigma$ & $(\alpha, \alpha, 2\alpha)$\\
$S(XS)$ & $(\frac{1}{2}+\alpha, 2\alpha, \frac{1}{2}+\alpha)$ & $S'(X'S')$ & $(2\alpha, \frac{1}{2}+\alpha, \frac{1}{2}+\alpha)$\\
$Z(XW)$ & $(\frac{1}{2}, \alpha, \frac{1}{2}+\alpha)$ & $Z'(X'W')$ & $(\alpha, \frac{1}{2}, \frac{1}{2}+\alpha)$\\
$Q(LQ)$ & $(\frac{1}{2}, \frac{1}{2}-\alpha, \frac{1}{2}+\alpha)$ & $Q'(LQ')$ & $(\frac{1}{2}-\alpha, \frac{1}{2}, \frac{1}{2}+\alpha)$\\
\bottomrule
\end{tabular}
\end{table}

我们必须对 (5.5.2) 所列 24 个空间群分别考虑 $\mathbf{G}^{\mathbf{k}_{B}}$ 的小表示的推导。把它们分为四组：

(a) $P2_{1}3\ (T^{4})$。

(b) $\{P4_{1}(C^{2}_{4}), P4_{3}(C^{4}_{4})\}$，$\{P3_{1}(C^{2}_{3}), P3_{2}(C^{3}_{3})\}$，$\{P3_{1}12(D^{3}_{3}), P3_{2}12(D^{5}_{3})\}$，$\{P3_{1}21(D^{4}_{3}), P3_{2}21(D^{6}_{3})\}$，$\{P6_{1}(C^{2}_{6}), P6_{5}(C^{3}_{6})\}$，$\{P6_{2}(C^{4}_{6}), P6_{4}(C^{5}_{6})\}$。

(c) $\{P4_{1}22(D^{3}_{4}), P4_{3}22(D^{7}_{4})\}$，$\{P4_{1}2_{1}2(D^{4}_{4}), P4_{3}2_{1}2(D^{8}_{4})\}$，$\{P6_{1}22(D^{2}_{6}), P6_{5}22(D^{3}_{6})\}$，$\{P6_{2}22(D^{4}_{6}), P6_{4}22(D^{5}_{6})\}$，$\{P4_{3}32(O^{6}), P4_{1}32(O^{7})\}$。

(d) $Pa3\ (T^{6}_{h})$。

其中同构对用大括号括起。(c)、(d) 中所列的群都不是任何更大空间群的不变子群。(a)、(b) 中所列的群虽然不是具有同形点群 $\mathbf{P}$ 的母群 $\mathbf{G}_{0}$ 的不变子群，却是某个（或某些）其他空间群 $\mathbf{G}'$ 的不变子群，$\mathbf{G}'$ 的同形点群的阶介于 $\mathbf{F}$ 的阶与 $\mathbf{P}$ 的阶之间。这些群 $\mathbf{G}'$ 在表 5.10 中指明。

\begin{table}[H]
\centering
\caption*{\textbf{表 5.10}\ 13 个空间群的不变性}
\begin{tabular}{@{}ll@{}}
\toprule
$\mathbf{G}$ & $\mathbf{G}'$\\
\midrule
$P2_{1}3\ (T^{4})$ & $P4_{3}32\ (O^{6}), P4_{1}32\ (O^{7}), Pa3\ (T^{6}_{h})$\\
$P4_{1}\ (C^{2}_{4})$ & $P4_{1}22\ (D^{3}_{4}), P4_{1}2_{1}2\ (D^{4}_{4})$\\
$P4_{3}\ (C^{4}_{4})$ & $P4_{3}22\ (D^{7}_{4}), P4_{3}2_{1}2\ (D^{8}_{4})$\\
$P6_{1}\ (C^{2}_{6}), P3_{1}12\ (D^{3}_{3}), P3_{1}21\ (D^{4}_{3}), P3_{1}\ (C^{2}_{3})$ & $P6_{1}22\ (D^{2}_{6})$\\
$P6_{5}\ (C^{3}_{6}), P3_{2}12\ (D^{5}_{3}), P3_{2}21\ (D^{6}_{3}), P3_{2}\ (C^{3}_{3})$ & $P6_{5}22\ (D^{3}_{6})$\\
$P6_{2}\ (C^{4}_{6}), P3_{2}12\ (D^{5}_{3}), P3_{2}21\ (D^{6}_{3}), P3_{2}\ (C^{3}_{3})$ & $P6_{2}22\ (D^{4}_{6})$\\
$P6_{4}\ (C^{5}_{6}), P3_{1}12\ (D^{3}_{3}), P3_{1}21\ (D^{4}_{3}), P3_{1}\ (C^{2}_{3})$ & $P6_{4}22\ (D^{5}_{6})$\\
\bottomrule
\end{tabular}
\end{table}

这一讨论结束了我们对表示域与基本域关系的处理；表 5.11 给出该信息的汇总（不同晶系的 $\Phi/\Omega$ 比值），供使用表 5.7 时参考。
\begin{figure}[H]
\centering
\includegraphics[width=0.86\textwidth]{c5_t11_p430.jpg}
\caption*{\textbf{表 5.11}\ 表示域与基本域的比值汇总（原书第 417 页扫描占位）。}
\end{figure}

